\section{Background and what is classical}\label{sec:background}

\paragraph{Engines.}
LLM engines batch per model iteration: running requests decode one token per iteration and waiting requests join between iterations~\citep{yu2022orca}; KV is stored in pages~\citep{kwon2023vllm}. SGLang keeps cached prefixes in a radix tree, evicts least-recently-used (LRU) leaves, and can order waiting requests by longest prefix match (LPM)~\citep{zheng2024sglang}. In SGLang 0.5.22 the default policy is FCFS and the default eviction rule LRU; LPM stops sorting its queue when more than 128 requests wait, a fallback that PEEK also documents, running its heavy-load cells at that depth~\citep{xie2026peek}; a waiting request does not lock its cached prefix, which is locked when the request is admitted; and under FCFS, including LPM above that threshold, an LRU refresh touches the prefixes of the first 128 waiting requests once every 32 scheduling passes (files under \fpath{python/sglang/srt/} at tag v0.5.22: \fpath{arg_groups/fields/schedule.py}, lines 82--98; \fpath{arg_groups/fields/memory.py}, lines 25--42; \fpath{managers/schedule_policy.py}, lines 133--135, 278--311 and 1591; \fpath{mem_cache/radix_cache.py}, lines 552--580; \fpath{environ.py}, line 594). vLLM 0.32.0 schedules FCFS, its default, or by priority, and a freed cached block keeps its hash until it is reallocated (\fpath{vllm/config/scheduler.py}, lines 21 and 159; \fpath{vllm/v1/core/block_pool.py}, lines 666--690 and 774--803; at tag v0.32.0). We first read these sources through a web tool that returns a language model's summary of each file (Appendix~\ref{app:prov}); the line references here and in Appendix~\ref{app:proofs} are to the files at the tags, read verbatim. Our synthesis of these rules is that in both engines a session's KV is protected only while one of its requests runs.

\paragraph{Thrashing and metastability.}
Denning showed that admitting one more program to a memory near saturation can collapse efficiency, and that keeping the working sets of the active programs within memory, which limits their number, prevents it~\citep{denning1968thrashing,denning1968working}. We call throughput as a function of the closed population $n$ the thrashing curve; the term is our framing. LRU does badly on loops larger than the cache~\citep{chou1985evaluation,smaragdakis1999eelru}. A metastable failure is a bad state that persists after its trigger is removed, sustained by a feedback loop under uncontrolled load; one of Bronson et al.'s case studies is a lost look-aside cache~\citep{bronson2021metastable}, and Huang et al.\ study such failures in production incidents and model them~\citep{huang2022metastable}. In their classification of sustaining effects, a trap in which the load itself lowers capacity is capacity-degradation amplification; one of their cases is garbage collection driven by queueing-induced memory pressure, and their experiments vary the trigger's magnitude against matched runs, the design we follow. In Bronson et al.'s terms, a system that serves its load but can be pushed by a trigger into the bad state is in a vulnerable state, and the bad state, kept by a sustaining effect after the trigger ends, is the metastable failure state. We also call the good state metastable, in the physics sense of a state that persists until a large enough perturbation. Gibbens, Hunt and Kelly showed hysteresis in loss networks with random alternative routing; they note that trunk reservation, which refuses alternatively routed calls on a link that is nearly full, removes it~\citep{gibbens1990bistability}. Denning measured the working-set window in a program's own virtual time, so that interruptions would not distort it~\citep{denning2005locality}; an LRU pool over idle KV ages its entries in real time, so the time a session waits in the queue counts against its residency.

\paragraph{What is classical, and what is specific to agent serving.}
The collapse past the working-set level, the remedy of limiting active jobs, and a bad state that outlives its trigger are classical, and so are metastability or bistability arising from a unimodal throughput curve under open load, hysteresis, and bounded admission as the cure, for random-access channels~\citep{kleinrock1975aloha,carleial1975bistable}, thrashing multiprogramming systems~\citep{courtois1975decomposability,wilkes1973dynamics,nelson1987stochastic,yokohira1989catastrophic} and service systems whose customers return~\citep{chan2014speedup,dong2015slowdowns,ingolfsson2020fluid}; in all of them the slowdown is a given function of occupancy. What is new on the theory side is Proposition~\ref{prop:A}, a bound specific to prefix caching, and the first-principles curve of \S\ref{sec:fp}, in which the slowdown is derived from the pool, the workload and the server, with its closed-form all-miss onset, speed law, fold criterion and trapped identity (Propositions~\ref{prop:nb}--\ref{prop:trapid}). The fluid picture of \S\ref{sec:fluid}, a unimodal throughput curve that gives one stable and one unstable equilibrium under open arrivals, is the classical fluid (drift) analysis of unstable random-access channels~\citep{kleinrock1975aloha,fayolle1977stability}; network congestion collapse, a stable bad state sustained by retransmissions, is its analogue~\citep{nagle1984congestion,jacobson1988congestion}. Replacing an open system's departure rate by its closed-system throughput is the classical decomposition (flow-equivalent) approximation, whose use to find unstable regimes of thrashing multiprogramming systems goes back to Courtois~\citep{courtois1975decomposability}, and choosing the number of admitted jobs from the throughput curve, by paging heuristics~\citep{denning1976optimal} or from measured throughput~\citep{heiss1991adaptive,schroeder2006mpl}, is multiprogramming-level load control. Two features of agent serving make these apply. The reusable KV is exactly the idle KV: a session in a tool call or waiting in the queue holds no running request, so memory-gated admission reclaims its cache. And reuse makes a hit much cheaper than a miss, so losing the cache multiplies the work of every later turn. Both features are the premise of earlier agent-serving systems~\citep{li2025continuum,abhyankar2024infercept,xia2026mori,bian2025tokencake,biswas2026sutradhara}. ThunderAgent describes agent collapse as working-set thrashing~\citep{kang2026thunderagent}, and NVIDIA Dynamo's port of its scheduler aims to avoid thrashing when memory cannot hold every agent's working set~\citep{dhanani2026dynamo}; EfficientAgent applies Denning's working-set principle and load control to the host KV tier of a fixed pool of agents and reports that waiting behind other agents stretches the gap between an agent's turns~\citep{shao2026efficientagent}. Hong et al.\ describe the feedback these features close, as a loop between eviction, cold prefill, queue length and session overlap, for RL rollouts under a concurrency ceiling, and interrupt it with a session cap~\citep{hong2026misat}; AgentKV, from its abstract, predicts the onset of collapse at high concurrency from available KV capacity over per-instance context demand~\citep{anon2026agentkv}. Under open Poisson arrivals on vLLM, KAIROS reports thrashing when its controller lowers the GPU clock (its Sec.~3.4 and Figs.~10--11) or when the arrival rate rises (its App.~A.4: 0.6 against 0.5 jobs/s), and avoids it by bounding the aggregate context of admitted agents (its Eq.~13: at most 0.95 of capacity, readmitting below 0.9)~\citep{yuan2026kairos}. In its Figs.~10--11, at 0.05 jobs/s, the run without thrashing avoidance stays in what KAIROS calls a persistent thrashing regime, completing 0.046 jobs/s, while a baseline at the maximum clock carries the same rate. Its SLO-aware boost (its Eq.~12) restores the maximum clock when an agent falls below its throughput target, but no frequency trace is shown, so whether the trigger was removed is not reported; the run has no burst trigger, no fixed clock, no replication and no matched arms. What we add is a controlled, replicated test at a fixed clock and pool, with a burst trigger, matched no-burst and capped arms and common random numbers (\S\ref{sec:e2}--\ref{sec:vllm}), the measured form the loop takes through the gap between a session's turns (\S\ref{sec:mechanism}; Proposition~\ref{prop:A} formalizes the last link in the large-population limit), and the first-principles curve that predicts where it closes (\S\ref{sec:fp}). We also compare published mechanisms in these terms in an earlier simulator (\S\ref{sec:mech}).
